% Anonymous example data file for template/resume-template.tex.
% Copy this file, replace every placeholder, then set:
% \newcommand{\ResumeDataFile}{template/your-data-file}
% % !TEX TS-program = xelatex
% !TEX encoding = UTF-8 Unicode
% !Mode:: "TeX:UTF-8"

\documentclass{resume}

\usepackage{zh_CN-systemfonts}
\usepackage{linespacing_fix}
\usepackage{cite}

\providecommand{\ResumeDataFile}{data/resume-data.example}
\input{\ResumeDataFile}

\begin{document}
\pagenumbering{gobble}
\RenderResume
\end{document}


\newcommand{\ResumeName}{张三}
\newcommand{\ResumePhone}{13800000000}
\newcommand{\ResumeEmail}{example@example.com}
\newcommand{\ResumeAge}{20岁}
\newcommand{\ResumeWebsite}{https://example.com}
\newcommand{\ResumeTitle}{某高校学生 / 研究人员}
\newcommand{\ResumePhoto}{images/example-profile.png}

\newcommand{\RenderResumeHeader}{%
  \resumeHeader
    {\ResumeName}
    {\ResumePhone}
    {\ResumeEmail}
    {\ResumeAge}
    {\ResumeWebsite}
    {\ResumeTitle}
    {\ResumePhoto}
}

\newcommand{\RenderEducationSection}{%
  \section{教育背景}
  \educationEntry{某某大学}{计算机科学与技术}{硕士}{2024.9 - 2027.6}
  \educationItemGap
  \educationEntry{某某大学}{软件工程}{本科}{2020.9 - 2024.6}
  \begin{itemize}
    \item \textbf{成绩}: 平均分 90.0 / 100，专业排名前 10\%
    \item \textbf{荣誉}: 校级奖学金；学科竞赛奖项；优秀学生干部
  \end{itemize}
}

\newcommand{\RenderResearchSection}{%
  \section{科研经历}
  \researchEntry{面向某类系统的高效计算优化研究}{某某实验室}{2024.1 - 至今}
  \begin{itemize}
    \researchTech{高性能计算、系统优化、并行计算、实验评测}
    \researchOutcomeFirst{完成相关论文投稿 / 发表，负责方法设计、实验实现与结果分析。}
    \researchWorkFirst{围绕目标系统中的计算瓶颈与资源利用问题，研究任务划分、数据流转、执行调度与性能稳定性之间的关系。}
    \researchWorkItem{建立系统化分析框架，从执行路径、缓存复用、并行度提升和资源约束四个方向评估优化策略的适用边界。}
    \researchWorkItem{构建真实负载驱动的实验验证链路，系统评估吞吐、延迟、资源占用和稳定性等指标。}
  \end{itemize}
}

\newcommand{\RenderProjectSection}{%
  \section{项目经历}
  \projectEntry{高性能计算系统原型实现}{https://example.com/project}
  \begin{itemize}
    \item 面向具体计算场景实现核心算法模块，优化数据访问路径、并行执行策略和端到端运行效率。
    \item 构建自动化测试与性能评测流程，完成多组输入规模下的正确性验证和性能对比。
  \end{itemize}
}

\newcommand{\RenderSkillsSection}{%
  \section{工作技能}
  \begin{itemize}
    \item \textbf{语言}: 英语 CET-6 / IELTS / TOEFL
    \item \textbf{工具与框架}: C / C++，Python，CUDA，Linux，Git
  \end{itemize}
}

\newcommand{\RenderResume}{%
  \RenderResumeHeader
  \RenderEducationSection
  \RenderResearchSection
  \RenderProjectSection
  \RenderSkillsSection
}
